\subsection*{ILOs - Experimental Causal Discovery}

We want the students to be able to:
\begin{itemize}
\item explain why correlation does not imply causation.
\item explain ~6 correlation inducing scenarios (e.g.\ cause-effect, effect-cause, feedback, 
    confounding, selection bias, deterministic constraints, etc.).
\item write down and explain the definition of causal effect (in real world terms).
\item explain how one can detect causal relationships experimentally.
\item explain randomized controlled trials and argue about advantages and disadvantages.
\item explain the difference between (conditional) observational distributions and interventional distributions, when they are different, when they are the same.
\item analyse real world studies and claims according to the concepts introduced above.

\end{itemize}


\subsection*{ILOs - Probability and Measure Theory}

We want the students to be able to:
\begin{itemize}
\item work with (elementary) probability theory, write down the main definitions and compute probabilities and expectation values.
\item translate measure theoretic notations, like measure integrals or expectation values, into familiar ones 
    in the discrete and absolute continuous cases.
\item understand, memorize and explain the 'core' definitions and constructions.
\item work with the Euclidean space $\R^D$ (as 'the' standard measurable space).
\item translate measure integrals to the two basic cases of either 1. discrete spaces/random variables/distributions (with probability mass functions) or 2. absolute continuous random variables (with probability densities w.r.t.\ Lebesgue measure) and compute those integrals (with the usual rules of anti-derivatives, etc.).
\item understand and apply the transformation rules for discrete and absolute continuous variables to compute mass functions/densities of push-forward measures (in the absolute continuous case only for differentiable bijections).
\end{itemize}


\subsection*{ILOs - Transition Probability Theory and Markov Kernels}

We want the students to be able to:
\begin{itemize}
\item explain what a Markov kernel is and how to represent them in the discrete and absolute continuous case.
\item give examples of Markov kernels.
\item properly work with discrete Markov kernels and linear Gaussians ones.
\item explain how Markov kernels relate to deterministic maps and probability distributions.
\item explain transition probability spaces and null-sets.
\item explain and work with the 5 constructions of Markov kernels from others 
    (composition, product, marginal, conditional, push-forward) 
    and how to work with them in the discrete and absolute continuous case.
\end{itemize}


\subsection*{ILOs - Conditional Independence}

We want the students to be able to:
\begin{itemize}
\item Write down the definition of conditional independence for conditional random variables 
    and work with them in the discrete and continuous cases.
\item Write down the separoid axioms for conditional independence of conditional random variables.
\item prove separoid axioms for random variables in the discrete and continuous case (in $\R^D$).
\item explain how one can get a given Markov kernel by a deterministic mapping and a probability distribution.
\item construct observational and interventional Markov kernels for SCMs.
\end{itemize}


\subsection*{ILOs - Graph Theory}

We want the students to be able to:
\begin{itemize}
\item write down the definition of CDMGs and explain the meaning.
\item explain all the different concepts of walks (directed, bidirected, collider, bifurcation, path, etc.).
\item explain all the different node/edge relationships (parents, children, etc.)
\item explain the notion of acyclicity
\item write down the definition of topological order and how it relates to acyclicity.
\item explain the different operations on graphs (hard/soft interventions, marginalization)
\item write down the definition of $id$-separation and check if sets are $id$-separated in given graphs.
\item write down the most important separoid axioms for $id$-separation, reason with them and prove them.
\end{itemize}


\subsection*{ILOs - Causal Bayesian Networks}

We want the students to be able to:
\begin{itemize}
\item write down the definition of causal Bayesian network with input and latent variables, motivate the meaning and reason about them.
\item work with the operations on causal Bayesian networks, like hard interventions, soft interventions, marginalization.
\item precisely state the global Markov property, write it down, reason about it and apply it to different causal Bayesian networks where the graph is given.
\item know the main ingredients and steps for the proof of the global Markov property (not necessarily every single step).
\item write down what interventional equivalence means and reason about it.
\item write down how a standard forms (SCM form, canonical form) of causal Bayesian network looks like and reason about it.
\end{itemize}


\subsection*{ILOs - Do-Calculus and Adjustment Formulas}

The student should be able to:
\begin{itemize}
\item write down the 3 do-calculus rules and reason with them.
\item prove the 3 rules of do-calculus using the global Markov property.
\item apply the 3 rules of do-calculus in an iterative manner on (small, not too complicated) graphs for the identification of causal effects (i.e. reducing the number $\doit$-terms).
\item state and apply the back-door covariate adjustment criterion and formula.
\item explain the general idea of the ID-algorithm.
\end{itemize}

\subsection*{Intended Learning Outcomes}

The student should be able to:
\begin{enumerate}
  \item Chapter 3
    \begin{itemize}
      \item know and understand the definition of $i\sigma$-separation;
      \item check if sets are $i\sigma$-separated in given graphs;
    \end{itemize}
  \item Chapter 6
    \begin{itemize}
    \item know and understand the definition of structural causal model (SCM);
    \item know the differences between endogenous variables, exogenous input variables and exogenous random variables;
    \item understand the implicit assumptions made when using a SCM;
    \item know and understand the definitions of (potential) outcomes, solutions, their Markov kernels, and solution functions of SCMs;
    \item know how to sample from a SCM;
    \item know and understand the definitions of the different variants of hard interventions on SCMs;
    \item understand how soft interventions can be modeled in SCMs;
    \item understand how intervention variables can be used to model different interventions in the same causal model;
    \item understand how to compose and decompose SCMs;
    \item propose an appropriate SCM for a given real-world system, and to appropriately model interventions on the system using the model;
    \item know and understand the definition and consequences of unique solvability for SCMs;
    \item know and understand the definition and consequences of an SCM being simple;
    \item be able to check whether a given SCM is uniquely solvable/simple;
    \item construct observational and interventional Markov kernels for simple SCMs;
    \item understand the different equivalence relations for SCMs (see also Chapter 8);
    \item know and understand the definition of marginalization of SCMs;
    \item construct a marginalized SCM;
    \item know and apply the basic commutation and compatibility properties of basic operations on SCMs (intervening, marginalization, preservation of unique solvability / simplicity / equivalence);
    \item know and understand the definition of the graph of an SCM;
    \item be able to draw the graph of a given SCM, and to construct an SCM with a given graph;
    \item know and apply the basic compatibility relations between the operation that maps an SCM to a graph and other basic SCM operations;
    \item give a causal interpretation of the graph of an SCM (see Chapter 9);
    \item know the definition of acyclic SCM;
    \item understand and appreciate the differences between different causal model classes.
  \end{itemize}
  \item Chapter 7
    \begin{itemize}
      \item know and understand the definitions of acyclification of a graph and of a simple SCM, and their relation;
      \item know the relation between $i\sigma$-separation and $id$-separation on an acyclification;
      \item know and understand the global Markov property for simple SCMs;
      \item apply the global Markov property for simple SCMs;
      \item know and understand the 3 do-calculus rules for simple SCMs;
      \item apply the 3 do-calculus rules for simple SCMs;
      \item state and apply the back-door covariate adjustment criterion and formula for simple SCMs.
    \end{itemize}
  \item Chapter 8
    \begin{itemize}
      \item construct a twin SCM;
      \item construct a twin graph;
      \item know and understand the basic commutation and compatibility properties of twinning with other basic operations on and properties of SCMs;
      \item be able to reason about counterfactuals by using the twin SCM;
      \item be aware of the empirical limitations of working with counterfactuals.
    \end{itemize}
  \item Chapter 9
    \begin{itemize}
      \item know and understand the definitions of indirect causal relations, direct causal relations, and common causes according to simple SCMs;
      \item understand the relation between having a common cause and being confounded;
      \item express the notion of ``no confounding'' in terms of potential outcomes;
      \item know how to detect the presence of such relations from Markov kernels induced by simple SCMs;
      \item deal formally with a situation in which only part of the variables in an SCM are observed, whereas the others are latent;
      \item formally model a randomized controlled trial with SCMs in various ways, and prove in the SCM formalism the main properties of RCTs;
      \item know and understand the definition of faithfulness;
      \item understand different ways of how faithfulness violations can occur;
      \item understand the LCD algorithm;
      \item prove the correctness of the LCD algorithm;
      \item understand the Y-structures approach to causal discovery;
      \item prove the correctness of the Y-structures approach to causal discovery;
      \item know how minimal separating/connecting sets relate to ancestral relations.
    \end{itemize}
  \item Chapter 10
    \begin{itemize}
      \item understand how to use the $G$ test for conditional independence;
      \item understand under which conditions the $G$ test is valid (at level $\alpha$);
      \item understand under which conditions the $G$ test is asymptotically consistent;
      \item implement the $G$ test in a computer program;
    \end{itemize}
  \item Chapter 11
    \begin{itemize}
      \item know and understand the definition of $\sigma$-inducing path;
      \item know and understand the definition of a DPAG;
      \item know which DMGs a DPAG represents;
      \item construct a DPAG representing a given DMG.
    \end{itemize}
\end{enumerate}

